\documentclass[11pt,a4paper]{article}
\usepackage[T1]{fontenc}
\usepackage{isabelle,isabellesym}

\usepackage{amsmath}
\usepackage{amssymb}
\usepackage{stmaryrd}

% this should be the last package used
\usepackage{pdfsetup}

\urlstyle{rm}
\isabellestyle{it}

\begin{document}

\title{Optional Sampling, Quadratic Variation and Transfer for Continuous-Time Martingales}
\author{Dmitriy Traytel}
\maketitle

\begin{abstract}
  What continuous-time martingale theory needs beyond the Archive's
  \texttt{Martingales} entry, on which this one builds.  The Archive's
  \texttt{Doob\_Convergence} proves the upcrossing inequality and the
  convergence theorem; the maximal inequality proved here is a different
  result, and optional sampling, quadratic variation and Vitali's theorem
  have no counterpart in either.

  Doob's maximal inequality is proved for discrete-time martingales, and with
  it a running-maximum bound usable as a dominating function.  Optional
  sampling is proved first for integrable discrete-time martingales, then for
  a continuous-time martingale sampled on a grid, then at a simple stopping
  time, and finally at an arbitrary bounded stopping time, by approximating
  the stopping time from above by its dyadic ceilings; the localised form,
  over an event of the past, comes along the way.  Optional stopping follows:
  the stopped process is again a martingale.  Exit times of closed sets by
  continuous adapted processes are shown to be stopping times.

  Quadratic variation of a square-integrable martingale is defined with its
  compensator, and the sampled process on a partition shown to be a
  square-integrable discrete-time martingale whose compensator is the sampled
  one --- which is what reduces a continuous-time statement to a discrete
  one.  Uniform integrability and Vitali's convergence theorem follow,
  with the practical criterion that a uniform higher-moment bound suffices,
  and the uniform integrability of a family of conditional expectations.

  A second half is algebraic.  The martingale property is closed under sums,
  differences, constant shifts, congruence up to almost-sure or eventual
  equality, time changes, coarser filtrations, and images under a bounded
  linear map, hence componentwise for vector- and matrix-valued processes;
  and it transports along image measures, products of two filtered measures,
  infinite products, restriction to a full-measure event, conditioning on an
  event of the past, and modification.  The natural filtration of a process,
  the semidirect product of a measure with a kernel, and a battery of
  integrability criteria complete the toolbox.  Karatzas and
  Shreve~\cite{KaratzasShreve} and Williams~\cite{Williams} are the textbook
  references.
\end{abstract}

\tableofcontents

\parindent 0pt\parskip 0.5ex

\input{session}

\bibliographystyle{abbrv}
\bibliography{root}

\end{document}
